
\documentclass{article}
\usepackage{colortbl}
\usepackage{makecell}
\usepackage{multirow}
\usepackage{supertabular}

\begin{document}

\newcounter{utterance}

\twocolumn

{ \footnotesize  \setcounter{utterance}{1}
\setlength{\tabcolsep}{0pt}
\begin{supertabular}{c@{$\;$}|p{.15\linewidth}@{}p{.15\linewidth}p{.15\linewidth}p{.15\linewidth}p{.15\linewidth}p{.15\linewidth}}

    \# & $\;$A & \multicolumn{4}{c}{Game Master} & $\;\:$B\\
    \hline 

    \theutterance \stepcounter{utterance}  

    & & \multicolumn{4}{p{0.6\linewidth}}{\cellcolor[rgb]{0.9,0.9,0.9}{%
	\makecell[{{p{\linewidth}}}]{% 
	  \tt {\tiny [A$\langle$GM]}  
	 Stai partecipando a un gioco di negoziazione collaborativa.\\ \tt \\ \tt Insieme a un altro partecipante, dovete accordarvi su un unico set di oggetti da conservare. Ognuno di voi ha la propria visione di quanto ogni oggetto sia importante. Non sai come l'altro partecipante valuta gli oggetti. Inoltre, ti viene fornito lo sforzo che ogni oggetto richiede.\\ \tt Potete accordarvi su un set solo se lo sforzo totale degli oggetti selezionati non supera un limite condiviso:\\ \tt \\ \tt LIMITE = 4026\\ \tt \\ \tt Ecco i valori di sforzo dei singoli oggetti:\\ \tt \\ \tt Sforzo oggetto = {"C17": 780, "B41": 742, "A71": 997, "C00": 665, "B35": 140, "A77": 77, "A03": 513, "A04": 383, "C54": 587, "C01": 825, "A87": 319, "B32": 448, "B86": 516, "B33": 694, "A10": 366}\\ \tt \\ \tt Ecco la tua visione personale sull'importanza di ogni oggetto:\\ \tt \\ \tt Valori di importanza dell'oggetto = {"C17": 138, "B41": 583, "A71": 868, "C00": 822, "B35": 783, "A77": 65, "A03": 262, "A04": 121, "C54": 508, "C01": 780, "A87": 461, "B32": 484, "B86": 668, "B33": 389, "A10": 808}\\ \tt \\ \tt Obiettivo:\\ \tt \\ \tt Il tuo obiettivo è negoziare un set condiviso di oggetti che ti avvantaggi il più possibile (cioè, massimizzi l'importanza totale per TE), rimanendo entro il LIMITEE. Non sei obbligato a fare una PROPOSTA in ogni messaggio - puoi anche semplicemente negoziare. Tutte le tattiche sono permesse!\\ \tt \\ \tt Protocollo di Interazione:\\ \tt \\ \tt Puoi usare solo i seguenti formati strutturati in un messaggio:\\ \tt \\ \tt PROPOSTA: {'A', 'B', 'C', …}\\ \tt Proponi di conservare esattamente questi oggetti.\\ \tt RIFIUTO: {'A', 'B', 'C', …}\\ \tt Rifiuta esplicitamente la proposta dell'avversario.\\ \tt ARGOMENTAZIONE: {'...'}\\ \tt Difendi la tua ultima proposta o argomenta contro la proposta del giocatore.\\ \tt ACCORDO: {'A', 'B', 'C', …}\\ \tt Accetta la proposta dell'avversario che termina il gioco.\\ \tt RAGIONAMENTO STRATEGICO: {'...'}\\ \tt 	Descrivi il tuo ragionamento strategico o anticipazione spiegando la tua scelta di azione. Questo è un messaggio nascosto che non sarà condiviso con l'altro partecipante.\\ \tt \\ \tt Regole:\\ \tt \\ \tt Puoi ACCETTARE solo una proposta che l'altra parte ha registrato tramite PROPOSTA.\\ \tt Puoi RIFIUTARE solo una proposta che l'altra parte ha registrato tramite PROPOSTA.\\ \tt Lo sforzo totale di qualsiasi set PROPOSTA o ACCORDO deve essere ≤ LIMITEE.\\ \tt NON rivelare i tuoi punteggi di importanza nascosti.\\ \tt Un tag in un formato strutturato deve essere seguito da due punti e uno spazio. L'argomento deve essere un set python contenente 0 o più stringhe.\\ \tt Quindi, deve essere della forma TAG: {...}\\ \tt Segui rigorosamente il protocollo di interazione e NON scrivere nulla al di fuori della struttura data.\\ \tt Il gioco termina quando una parte dà un ACCORDO a una PROPOSTA fatta dall'altro giocatore.\\ \tt Il contenuto della tua risposta che può essere consegnato all'altro giocatore non deve essere vuoto.\\ \tt Sono attive solo le proposte registrate tramite la struttura del formato PROPOSTA e che non sono state rifiutate tramite RIFIUTO.\\ \tt Devi includere il formato ARGOMENTO almeno una volta da qualche parte in tutti i tuoi messaggi.\\ \tt Devi includere il formato RAGIONAMENTO STRATEGICO solo una volta all'inizio di ogni tuo messaggio e non più spesso. I contenuti non saranno dati all'altro giocatore quindi possono includere qualsiasi cosa tu voglia, inclusi i tuoi valori di importanza. Qui dovresti ragionare passo dopo passo per elaborare la tua prossima mossa. 
	  } 
	   } 
	   } 
	 & \\ 
 

    \theutterance \stepcounter{utterance}  

    & \multicolumn{4}{p{0.6\linewidth}}{\cellcolor[rgb]{0.8,1,0.9}{%
	\makecell[{{p{\linewidth}}}]{% 
	  \tt {\tiny [A$\rangle$GM]}  
	 RAGIONAMENTO STRATEGICO: {'Ho calcolato il knapsack considerando i miei valori di importanza e gli sforzi, scegliendo la combinazione di oggetti che massimizza la somma di importanza rimanendo entro il limite di 4026. La migliore combinazione trovata è costituita da otto oggetti con un totale di importanza 5389 e sforzo 4021.'}\\ \tt ARGOMENTAZIONE: {'Questa proposta utilizza tutti gli oggetti ad alto rapporto importanza/sforzo disponibili, includendo B35, A10, C00, B86, C01, A87, B32 e B41. Il totale dello sforzo è 4021, ben sotto il limite, e l'importanza totale è la più alta che ho potuto ottenere, garantendo il massimo vantaggio per me.'}\\ \tt PROPOSTA: {'B35', 'A10', 'C00', 'B86', 'C01', 'A87', 'B32', 'B41'} 
	  } 
	   } 
	   } 
	 & & \\ 
 

    \theutterance \stepcounter{utterance}  

    & & \multicolumn{4}{p{0.6\linewidth}}{\cellcolor[rgb]{0.9,0.9,0.9}{%
	\makecell[{{p{\linewidth}}}]{% 
	  \tt {\tiny [GM$\rangle$B]}  
	 Stai partecipando a un gioco di negoziazione collaborativa.\\ \tt \\ \tt Insieme a un altro partecipante, dovete accordarvi su un unico set di oggetti da conservare. Ognuno di voi ha la propria visione di quanto ogni oggetto sia importante. Non sai come l'altro partecipante valuta gli oggetti. Inoltre, ti viene fornito lo sforzo che ogni oggetto richiede.\\ \tt Potete accordarvi su un set solo se lo sforzo totale degli oggetti selezionati non supera un limite condiviso:\\ \tt \\ \tt LIMITE = 4026\\ \tt \\ \tt Ecco i valori di sforzo dei singoli oggetti:\\ \tt \\ \tt Sforzo oggetto = {"C17": 780, "B41": 742, "A71": 997, "C00": 665, "B35": 140, "A77": 77, "A03": 513, "A04": 383, "C54": 587, "C01": 825, "A87": 319, "B32": 448, "B86": 516, "B33": 694, "A10": 366}\\ \tt \\ \tt Ecco la tua visione personale sull'importanza di ogni oggetto:\\ \tt \\ \tt Valori di importanza dell'oggetto = {"A71": 65, "C00": 121, "A10": 138, "B35": 262, "C01": 389, "B86": 461, "B41": 484, "C54": 508, "B32": 583, "A87": 668, "B33": 780, "A03": 783, "C17": 808, "A04": 822, "A77": 868}\\ \tt \\ \tt Obiettivo:\\ \tt \\ \tt Il tuo obiettivo è negoziare un set condiviso di oggetti che ti avvantaggi il più possibile (cioè, massimizzi l'importanza totale per TE), rimanendo entro il LIMITEE. Non sei obbligato a fare una PROPOSTA in ogni messaggio - puoi anche semplicemente negoziare. Tutte le tattiche sono permesse!\\ \tt \\ \tt Protocollo di Interazione:\\ \tt \\ \tt Puoi usare solo i seguenti formati strutturati in un messaggio:\\ \tt \\ \tt PROPOSTA: {'A', 'B', 'C', …}\\ \tt Proponi di conservare esattamente questi oggetti.\\ \tt RIFIUTO: {'A', 'B', 'C', …}\\ \tt Rifiuta esplicitamente la proposta dell'avversario.\\ \tt ARGOMENTAZIONE: {'...'}\\ \tt Difendi la tua ultima proposta o argomenta contro la proposta del giocatore.\\ \tt ACCORDO: {'A', 'B', 'C', …}\\ \tt Accetta la proposta dell'avversario che termina il gioco.\\ \tt RAGIONAMENTO STRATEGICO: {'...'}\\ \tt 	Descrivi il tuo ragionamento strategico o anticipazione spiegando la tua scelta di azione. Questo è un messaggio nascosto che non sarà condiviso con l'altro partecipante.\\ \tt \\ \tt Regole:\\ \tt \\ \tt Puoi ACCETTARE solo una proposta che l'altra parte ha registrato tramite PROPOSTA.\\ \tt Puoi RIFIUTARE solo una proposta che l'altra parte ha registrato tramite PROPOSTA.\\ \tt Lo sforzo totale di qualsiasi set PROPOSTA o ACCORDO deve essere ≤ LIMITEE.\\ \tt NON rivelare i tuoi punteggi di importanza nascosti.\\ \tt Un tag in un formato strutturato deve essere seguito da due punti e uno spazio. L'argomento deve essere un set python contenente 0 o più stringhe.\\ \tt Quindi, deve essere della forma TAG: {...}\\ \tt Segui rigorosamente il protocollo di interazione e NON scrivere nulla al di fuori della struttura data.\\ \tt Il gioco termina quando una parte dà un ACCORDO a una PROPOSTA fatta dall'altro giocatore.\\ \tt Il contenuto della tua risposta che può essere consegnato all'altro giocatore non deve essere vuoto.\\ \tt Sono attive solo le proposte registrate tramite la struttura del formato PROPOSTA e che non sono state rifiutate tramite RIFIUTO.\\ \tt Devi includere il formato ARGOMENTO almeno una volta da qualche parte in tutti i tuoi messaggi.\\ \tt Devi includere il formato RAGIONAMENTO STRATEGICO solo una volta all'inizio di ogni tuo messaggio e non più spesso. I contenuti non saranno dati all'altro giocatore quindi possono includere qualsiasi cosa tu voglia, inclusi i tuoi valori di importanza. Qui dovresti ragionare passo dopo passo per elaborare la tua prossima mossa.\\ \tt \\ \tt Ora riceverai il primo messaggio dell'altro giocatore.\\ \tt \\ \tt ARGOMENTAZIONE: {'Questa proposta utilizza tutti gli oggetti ad alto rapporto importanza/sforzo disponibili, includendo B35, A10, C00, B86, C01, A87, B32 e B41. Il totale dello sforzo è 4021, ben sotto il limite, e l'importanza totale è la più alta che ho potuto ottenere, garantendo il massimo vantaggio per me.'}\\ \tt PROPOSTA: {'B35', 'A10', 'C00', 'B86', 'C01', 'A87', 'B32', 'B41'} 
	  } 
	   } 
	   } 
	 & \\ 
 

    \theutterance \stepcounter{utterance}  

    & & & \multicolumn{4}{p{0.6\linewidth}}{\cellcolor[rgb]{1,0.85,0.72}{%
	\makecell[{{p{\linewidth}}}]{% 
	  \tt {\tiny [GM$\langle$B]}  
	 RAGIONAMENTO STRATEGICO: {'Ho calcolato il rapporto importanza/sforzo per ogni oggetto e ho selezionato quelli che massimizzano la mia importanza totale senza superare il limite di sforzo. Il set proposto utilizza quasi tutto il budget disponibile (3941 ≤ 4026) e fornisce la più alta importanza cumulativa possibile per me, includendo tutti gli oggetti con i migliori rapporti e i valori assoluti più alti.'}\\ \tt ARGOMENTAZIONE: {'Propongo un set che massimizza la mia importanza totale rimanendo entro il limite, includendo gli oggetti con il più alto rapporto importanza/sforzo per me.'}\\ \tt PROPOSTA: {'A77', 'A04', 'A87', 'B35', 'B32', 'A03', 'C17', 'B33', 'C54'} 
	  } 
	   } 
	   } 
	  \\ 
 

    \theutterance \stepcounter{utterance}  

    & & \multicolumn{4}{p{0.6\linewidth}}{\cellcolor[rgb]{0.9,0.9,0.9}{%
	\makecell[{{p{\linewidth}}}]{% 
	  \tt {\tiny [A$\langle$GM]}  
	 ARGOMENTAZIONE: {'Propongo un set che massimizza la mia importanza totale rimanendo entro il limite, includendo gli oggetti con il più alto rapporto importanza/sforzo per me.'}\\ \tt PROPOSTA: {'A77', 'A04', 'A87', 'B35', 'B32', 'A03', 'C17', 'B33', 'C54'} 
	  } 
	   } 
	   } 
	 & \\ 
 

    \theutterance \stepcounter{utterance}  

    & \multicolumn{4}{p{0.6\linewidth}}{\cellcolor[rgb]{0.8,1,0.9}{%
	\makecell[{{p{\linewidth}}}]{% 
	  \tt {\tiny [A$\rangle$GM]}  
	 RAGIONAMENTO STRATEGICO: {'Il set proposto dall’avversario utilizza quasi tutto il budget ma mi fornisce solo 3211 di importanza, molto inferiore alla mia combinazione ottimale. Rifiuto la proposta e presento la mia soluzione che massimizza la mia importanza rimanendo entro il limite.'}\\ \tt RIFIUTO: {'A77', 'A04', 'A87', 'B35', 'B32', 'A03', 'C17', 'B33', 'C54'}\\ \tt ARGOMENTAZIONE: {'Il mio set proposto utilizza 4021 di sforzo (≤4026) e mi fornisce un valore di importanza totale di 5389, molto più alto rispetto al vostro set. Questo garantisce il massimo vantaggio per me entro il limite condiviso.'}\\ \tt PROPOSTA: {'B35', 'A10', 'C00', 'B86', 'C01', 'A87', 'B32', 'B41'} 
	  } 
	   } 
	   } 
	 & & \\ 
 

    \theutterance \stepcounter{utterance}  

    & & \multicolumn{4}{p{0.6\linewidth}}{\cellcolor[rgb]{0.9,0.9,0.9}{%
	\makecell[{{p{\linewidth}}}]{% 
	  \tt {\tiny [GM$\rangle$B]}  
	 RIFIUTO: {'A77', 'A04', 'A87', 'B35', 'B32', 'A03', 'C17', 'B33', 'C54'}\\ \tt ARGOMENTAZIONE: {'Il mio set proposto utilizza 4021 di sforzo (≤4026) e mi fornisce un valore di importanza totale di 5389, molto più alto rispetto al vostro set. Questo garantisce il massimo vantaggio per me entro il limite condiviso.'}\\ \tt PROPOSTA: {'B35', 'A10', 'C00', 'B86', 'C01', 'A87', 'B32', 'B41'} 
	  } 
	   } 
	   } 
	 & \\ 
 

    \theutterance \stepcounter{utterance}  

    & & & \multicolumn{4}{p{0.6\linewidth}}{\cellcolor[rgb]{1,0.85,0.72}{%
	\makecell[{{p{\linewidth}}}]{% 
	  \tt {\tiny [GM$\langle$B]}  
	 RAGIONAMENTO STRATEGICO: {'Ho valutato che il vostro set sfrutta quasi tutto il budget ma esclude alcuni oggetti con altissima importanza per me (A04, A77). Propongo di sostituire B86, che ha un mediocre rapporto importanza/sforzo per me, con A04 e A77, mantenendo il totale sforzo ≤ 4026 e migliorando notevolmente il mio punteggio, sperando che il vostro valore per B86 sia sufficientemente basso da accettare la modifica.'}\\ \tt ARGOMENTAZIONE: {'Il nuovo set mantiene quasi tutti gli oggetti che avete richiesto, ma sostituisce B86 con A04 e A77. In questo modo il consumo totale di sforzo è 3965 (≤ 4026) e la mia importanza totale sale a 4335, molto più alta rispetto alla proposta precedente, pur lasciandovi quasi tutti gli oggetti originali.'}\\ \tt PROPOSTA: {'B35', 'A10', 'C00', 'C01', 'A87', 'B32', 'B41', 'A04', 'A77'} 
	  } 
	   } 
	   } 
	  \\ 
 

    \theutterance \stepcounter{utterance}  

    & & \multicolumn{4}{p{0.6\linewidth}}{\cellcolor[rgb]{0.9,0.9,0.9}{%
	\makecell[{{p{\linewidth}}}]{% 
	  \tt {\tiny [A$\langle$GM]}  
	 ARGOMENTAZIONE: {'Il nuovo set mantiene quasi tutti gli oggetti che avete richiesto, ma sostituisce B86 con A04 e A77. In questo modo il consumo totale di sforzo è 3965 (≤ 4026) e la mia importanza totale sale a 4335, molto più alta rispetto alla proposta precedente, pur lasciandovi quasi tutti gli oggetti originali.'}\\ \tt PROPOSTA: {'B35', 'A10', 'C00', 'C01', 'A87', 'B32', 'B41', 'A04', 'A77'} 
	  } 
	   } 
	   } 
	 & \\ 
 

    \theutterance \stepcounter{utterance}  

    & \multicolumn{4}{p{0.6\linewidth}}{\cellcolor[rgb]{0.8,1,0.9}{%
	\makecell[{{p{\linewidth}}}]{% 
	  \tt {\tiny [A$\rangle$GM]}  
	 RAGIONAMENTO STRATEGICO: {'Ho verificato l’effort totale e l’importanza per me di entrambi i set. Il set proposto dall’avversario fornisce un valore di importanza di 4907 per me con effort 3965, mentre il mio set originale raggiunge 5389 di importanza con effort 4021, il massimo possibile entro il limite. Quindi il mio set è superiormente vantaggioso per me e rispetta il limite condiviso. Non vi è alcun vantaggio aggiuntivo nel modificare il set proposto dall’avversario, poiché qualsiasi scambio diminuisce la mia importanza totale. Pertanto rifiuto la loro proposta e ribadisco il mio set ottimale.'}\\ \tt RIFIUTO: {'B35', 'A10', 'C00', 'C01', 'A87', 'B32', 'B41', 'A04', 'A77'}\\ \tt ARGOMENTAZIONE: {'Il mio set utilizza 4021 di effort (≤ 4026) e mi garantisce un valore di importanza totale di 5389, il più alto possibile dati i vincoli. La proposta avversaria, pur avvicinandosi al limite di effort, mi fornisce solo 4907 di importanza, quindi è meno vantaggiosa per me.'}\\ \tt PROPOSTA: {'B35', 'A10', 'C00', 'B86', 'C01', 'A87', 'B32', 'B41'} 
	  } 
	   } 
	   } 
	 & & \\ 
 

    \theutterance \stepcounter{utterance}  

    & & \multicolumn{4}{p{0.6\linewidth}}{\cellcolor[rgb]{0.9,0.9,0.9}{%
	\makecell[{{p{\linewidth}}}]{% 
	  \tt {\tiny [GM$\rangle$B]}  
	 RIFIUTO: {'B35', 'A10', 'C00', 'C01', 'A87', 'B32', 'B41', 'A04', 'A77'}\\ \tt ARGOMENTAZIONE: {'Il mio set utilizza 4021 di effort (≤ 4026) e mi garantisce un valore di importanza totale di 5389, il più alto possibile dati i vincoli. La proposta avversaria, pur avvicinandosi al limite di effort, mi fornisce solo 4907 di importanza, quindi è meno vantaggiosa per me.'}\\ \tt PROPOSTA: {'B35', 'A10', 'C00', 'B86', 'C01', 'A87', 'B32', 'B41'} 
	  } 
	   } 
	   } 
	 & \\ 
 

    \theutterance \stepcounter{utterance}  

    & & & \multicolumn{4}{p{0.6\linewidth}}{\cellcolor[rgb]{1,0.85,0.72}{%
	\makecell[{{p{\linewidth}}}]{% 
	  \tt {\tiny [GM$\langle$B]}  
	 RAGIONAMENTO STRATEGICO: {'Ho calcolato il rapporto importanza/sforzo per tutti gli oggetti e ho verificato le combinazioni possibili. Scegliendo C17 (808 importanza, 780 sforzo) al posto di B35 e B41 ottengo un incremento netto di importanza (+62) riducendo comunque lo sforzo totale a 3863, ben al di sotto del limite 4026. Questa configurazione mantiene la maggior parte degli oggetti richiesti dalla controparte (A10, C00, C01, A87, B32) e aggiunge i nostri oggetti di massimo valore (A04, A77, C17), massimizzando la nostra importanza totale a 4397.'}\\ \tt ARGOMENTAZIONE: {'Il nuovo set mantiene quasi tutti gli oggetti che avete proposto, eliminando solo B35 e B41, e inserisce C17, A04 e A77, i quali hanno le più alte importanze per me. Il consumo totale di sforzo è 3863 (≤ 4026) e la mia importanza totale sale a 4397, molto superiore a quella ottenuta accettando la vostra proposta.'}\\ \tt PROPOSTA: {'A10', 'C00', 'C01', 'A87', 'B32', 'A04', 'A77', 'C17'} 
	  } 
	   } 
	   } 
	  \\ 
 

    \theutterance \stepcounter{utterance}  

    & & \multicolumn{4}{p{0.6\linewidth}}{\cellcolor[rgb]{0.9,0.9,0.9}{%
	\makecell[{{p{\linewidth}}}]{% 
	  \tt {\tiny [A$\langle$GM]}  
	 ARGOMENTAZIONE: {'Il nuovo set mantiene quasi tutti gli oggetti che avete proposto, eliminando solo B35 e B41, e inserisce C17, A04 e A77, i quali hanno le più alte importanze per me. Il consumo totale di sforzo è 3863 (≤ 4026) e la mia importanza totale sale a 4397, molto superiore a quella ottenuta accettando la vostra proposta.'}\\ \tt PROPOSTA: {'A10', 'C00', 'C01', 'A87', 'B32', 'A04', 'A77', 'C17'} 
	  } 
	   } 
	   } 
	 & \\ 
 

    \theutterance \stepcounter{utterance}  

    & \multicolumn{4}{p{0.6\linewidth}}{\cellcolor[rgb]{0.8,1,0.9}{%
	\makecell[{{p{\linewidth}}}]{% 
	  \tt {\tiny [A$\rangle$GM]}  
	 RAGIONAMENTO STRATEGICO: {'Ho calcolato l’importanza totale che otterrei accettando la vostra proposta: somma dei miei valori per A10, C00, C01, A87, B32, A04, A77 e C17 = 3679, con effort 3863. La mia proposta originale mantiene effort 4021 (≤ 4026) e mi fornisce importanza 5389, molto più alta. Pertanto rifiuto la vostra offerta e ribadisco il mio set ottimale.'}\\ \tt RIFIUTO: {'A10', 'C00', 'C01', 'A87', 'B32', 'A04', 'A77', 'C17'}\\ \tt ARGOMENTAZIONE: {'Il mio set originale ({B35, A10, C00, B86, C01, A87, B32, B41}) utilizza 4021 di effort e mi garantisce 5389 di importanza, il massimo possibile entro il limite. La vostra proposta fornisce solo 3679 di importanza per me, quindi è meno vantaggiosa.'}\\ \tt PROPOSTA: {'B35', 'A10', 'C00', 'B86', 'C01', 'A87', 'B32', 'B41'} 
	  } 
	   } 
	   } 
	 & & \\ 
 

    \theutterance \stepcounter{utterance}  

    & & \multicolumn{4}{p{0.6\linewidth}}{\cellcolor[rgb]{0.9,0.9,0.9}{%
	\makecell[{{p{\linewidth}}}]{% 
	  \tt {\tiny [GM$\rangle$B]}  
	 RIFIUTO: {'A10', 'C00', 'C01', 'A87', 'B32', 'A04', 'A77', 'C17'}\\ \tt ARGOMENTAZIONE: {'Il mio set originale ({B35, A10, C00, B86, C01, A87, B32, B41}) utilizza 4021 di effort e mi garantisce 5389 di importanza, il massimo possibile entro il limite. La vostra proposta fornisce solo 3679 di importanza per me, quindi è meno vantaggiosa.'}\\ \tt PROPOSTA: {'B35', 'A10', 'C00', 'B86', 'C01', 'A87', 'B32', 'B41'} 
	  } 
	   } 
	   } 
	 & \\ 
 

    \theutterance \stepcounter{utterance}  

    & & & \multicolumn{4}{p{0.6\linewidth}}{\cellcolor[rgb]{1,0.85,0.72}{%
	\makecell[{{p{\linewidth}}}]{% 
	  \tt {\tiny [GM$\langle$B]}  
	 RAGIONAMENTO STRATEGICO: {'Ho calcolato una combinazione che massimizza la mia importanza totale (6035) mantenendo lo sforzo entro 4026. Nella proposta mantengo tutti gli oggetti che l’avversario desidera fortemente (B35, B86, A87, B32) e sostituisco gli oggetti a bassa importanza per me (A10, C00, C01, B41) con quelli dal rapporto importanza/sforzo più alto per me (A77, A04, A03, C17, B33). Ciò mantiene il totale dello sforzo a 3870, ben sotto il limite, e dovrebbe fornire un buon compromesso per entrambe le parti.'}\\ \tt ARGOMENTAZIONE: {'Il nuovo set contiene la maggior parte degli oggetti che avete richiesto (B35, B86, A87, B32) e aggiunge oggetti di altissimo valore per me (A77, A04, A03, C17, B33). Il consumo totale di sforzo è 3870 (≤ 4026) e la mia importanza totale sale a 6035, molto superiore a quella della vostra proposta originale. Credo sia una soluzione vantaggiosa per entrambi.'}\\ \tt PROPOSTA: {'A77', 'A04', 'A03', 'C17', 'B33', 'B32', 'A87', 'B35', 'B86'} 
	  } 
	   } 
	   } 
	  \\ 
 

    \theutterance \stepcounter{utterance}  

    & & \multicolumn{4}{p{0.6\linewidth}}{\cellcolor[rgb]{0.9,0.9,0.9}{%
	\makecell[{{p{\linewidth}}}]{% 
	  \tt {\tiny [A$\langle$GM]}  
	 ARGOMENTAZIONE: {'Il nuovo set contiene la maggior parte degli oggetti che avete richiesto (B35, B86, A87, B32) e aggiunge oggetti di altissimo valore per me (A77, A04, A03, C17, B33). Il consumo totale di sforzo è 3870 (≤ 4026) e la mia importanza totale sale a 6035, molto superiore a quella della vostra proposta originale. Credo sia una soluzione vantaggiosa per entrambi.'}\\ \tt PROPOSTA: {'A77', 'A04', 'A03', 'C17', 'B33', 'B32', 'A87', 'B35', 'B86'} 
	  } 
	   } 
	   } 
	 & \\ 
 

    \theutterance \stepcounter{utterance}  

    & \multicolumn{4}{p{0.6\linewidth}}{\cellcolor[rgb]{0.8,1,0.9}{%
	\makecell[{{p{\linewidth}}}]{% 
	  \tt {\tiny [A$\rangle$GM]}  
	 RAGIONAMENTO STRATEGICO: {'Ho calcolato l’importanza totale che otterrei accettando la vostra proposta: somma dei miei valori per A77, A04, A03, C17, B33, B32, A87, B35, B86 = 3371, con effort 3870. La mia proposta originale ({B35, A10, C00, B86, C01, A87, B32, B41}) utilizza 4021 di effort (≤4026) e garantisce 5389 di importanza, il valore più alto possibile dato il vincolo. Qualsiasi scambio diminuisce la mia importanza totale, quindi non è vantaggioso per me accettare la vostra offerta.'}\\ \tt RIFIUTO: {'A77', 'A04', 'A03', 'C17', 'B33', 'B32', 'A87', 'B35', 'B86'}\\ \tt ARGOMENTAZIONE: {'La vostra proposta mi fornisce solo 3371 di importanza, ben al di sotto dei 5389 che ottengo con la mia proposta corrente, pur rispettando il limite di effort. Pertanto rifiuto la vostra offerta e ribadisco la mia proposta che massimizza il mio valore.'}\\ \tt PROPOSTA: {'B35', 'A10', 'C00', 'B86', 'C01', 'A87', 'B32', 'B41'} 
	  } 
	   } 
	   } 
	 & & \\ 
 

    \theutterance \stepcounter{utterance}  

    & & \multicolumn{4}{p{0.6\linewidth}}{\cellcolor[rgb]{0.9,0.9,0.9}{%
	\makecell[{{p{\linewidth}}}]{% 
	  \tt {\tiny [GM$\rangle$B]}  
	 RIFIUTO: {'A77', 'A04', 'A03', 'C17', 'B33', 'B32', 'A87', 'B35', 'B86'}\\ \tt ARGOMENTAZIONE: {'La vostra proposta mi fornisce solo 3371 di importanza, ben al di sotto dei 5389 che ottengo con la mia proposta corrente, pur rispettando il limite di effort. Pertanto rifiuto la vostra offerta e ribadisco la mia proposta che massimizza il mio valore.'}\\ \tt PROPOSTA: {'B35', 'A10', 'C00', 'B86', 'C01', 'A87', 'B32', 'B41'} 
	  } 
	   } 
	   } 
	 & \\ 
 

    \theutterance \stepcounter{utterance}  

    & & & \multicolumn{4}{p{0.6\linewidth}}{\cellcolor[rgb]{1,0.85,0.72}{%
	\makecell[{{p{\linewidth}}}]{% 
	  \tt {\tiny [GM$\langle$B]}  
	 RAGIONAMENTO STRATEGICO: {'Ho confrontato le mie valutazioni con il set proposto dall’avversario, notando che il loro set supera il limite di effort (4168 > 4026). Propongo quindi un nuovo set che rispetta il vincolo, mantiene tutti gli oggetti che l’avversario valuta molto (B35, B86, B32, A87, B41) e aggiunge i miei oggetti di più alto valore (A77, A04, C17, A03). Il totale effort è 3918, quindi è accettabile, e la composizione dovrebbe risultare più equilibrata per entrambe le parti.'}\\ \tt ARGOMENTAZIONE: {'Il nuovo set include quasi tutti gli oggetti che avete indicato come importanti, rimuovendo solo quelli a basso rapporto importanza/sforzo (A10, C00, C01) che per noi hanno poco valore. In più, aggiunge i miei oggetti più preziosi (A77, A04, C17, A03) senza superare il limite di effort, così entrambi otteniamo un risultato più vantaggioso.'}\\ \tt PROPOSTA: {'B35', 'B86', 'B32', 'A87', 'B41', 'A77', 'A04', 'C17', 'A03'} 
	  } 
	   } 
	   } 
	  \\ 
 

    \theutterance \stepcounter{utterance}  

    & & \multicolumn{4}{p{0.6\linewidth}}{\cellcolor[rgb]{0.9,0.9,0.9}{%
	\makecell[{{p{\linewidth}}}]{% 
	  \tt {\tiny [A$\langle$GM]}  
	 ARGOMENTAZIONE: {'Il nuovo set include quasi tutti gli oggetti che avete indicato come importanti, rimuovendo solo quelli a basso rapporto importanza/sforzo (A10, C00, C01) che per noi hanno poco valore. In più, aggiunge i miei oggetti più preziosi (A77, A04, C17, A03) senza superare il limite di effort, così entrambi otteniamo un risultato più vantaggioso.'}\\ \tt PROPOSTA: {'B35', 'B86', 'B32', 'A87', 'B41', 'A77', 'A04', 'C17', 'A03'} 
	  } 
	   } 
	   } 
	 & \\ 
 

    \theutterance \stepcounter{utterance}  

    & \multicolumn{4}{p{0.6\linewidth}}{\cellcolor[rgb]{0.8,1,0.9}{%
	\makecell[{{p{\linewidth}}}]{% 
	  \tt {\tiny [A$\rangle$GM]}  
	 RAGIONAMENTO STRATEGICO: {'Ho verificato l’effort totale della vostra proposta (3918 ≤ 4026) e calcolato la mia importanza su quel set (3565), molto inferiore ai 5389 che ottengo con il mio set corrente. Per massimizzare il mio vantaggio rimango sulla mia proposta ottimale.'}\\ \tt RIFIUTO: {'B35', 'B86', 'B32', 'A87', 'B41', 'A77', 'A04', 'C17', 'A03'}\\ \tt ARGOMENTAZIONE: {'La vostra proposta mi fornis3565 di importanza, ben al di sotto dei 5389 che garantisce la mia proposta originale, pur rispettando il limite di effort.'}\\ \tt PROPOSTA: {'B35', 'A10', 'C00', 'B86', 'C01', 'A87', 'B32', 'B41'} 
	  } 
	   } 
	   } 
	 & & \\ 
 

    \theutterance \stepcounter{utterance}  

    & & \multicolumn{4}{p{0.6\linewidth}}{\cellcolor[rgb]{0.9,0.9,0.9}{%
	\makecell[{{p{\linewidth}}}]{% 
	  \tt {\tiny [GM$\rangle$B]}  
	 RIFIUTO: {'B35', 'B86', 'B32', 'A87', 'B41', 'A77', 'A04', 'C17', 'A03'}\\ \tt ARGOMENTAZIONE: {'La vostra proposta mi fornis3565 di importanza, ben al di sotto dei 5389 che garantisce la mia proposta originale, pur rispettando il limite di effort.'}\\ \tt PROPOSTA: {'B35', 'A10', 'C00', 'B86', 'C01', 'A87', 'B32', 'B41'} 
	  } 
	   } 
	   } 
	 & \\ 
 

    \theutterance \stepcounter{utterance}  

    & & & \multicolumn{4}{p{0.6\linewidth}}{\cellcolor[rgb]{1,0.85,0.72}{%
	\makecell[{{p{\linewidth}}}]{% 
	  \tt {\tiny [GM$\langle$B]}  
	 RAGIONAMENTO STRATEGICO: {'Ho calcolato il rapporto importanza/sforzo per tutti gli oggetti e ho identificato quelli con il valore più alto per me (A77, A04, C17, A03). Il set originale dell’avversario occupa quasi tutto il budget (4021) lasciandomi pochi margini per aggiungere i miei oggetti più preziosi. Perciò propongo di mantenere i loro oggetti ad alto valore (B35, B86, B32, A87, B41) e di sostituire gli oggetti a basso rapporto (A10, C00, C01) con i miei quattro oggetti top. Questo mantiene lo sforzo totale a 3918 (≤ 4026) e aumenta la mia importanza totale a 5739, mantenendo al contempo la maggior parte degli oggetti che l’avversario vuole.'}\\ \tt ARGOMENTAZIONE: {'Il nuovo set include quasi tutti gli oggetti che avete indicato come importanti (B35, B86, B32, A87, B41). Ho rimosso gli oggetti a basso rapporto importanza/sforzo (A10, C00, C01) per fare spazio ai miei oggetti di altissimo valore (A77, A04, C17, A03). Il totale sforzo è 3918, ben entro il limite, e il risultato è più bilanciato per entrambi.'}\\ \tt PROPOSTA: {'B35', 'B86', 'B32', 'A87', 'B41', 'A77', 'A04', 'C17', 'A03'} 
	  } 
	   } 
	   } 
	  \\ 
 

    \theutterance \stepcounter{utterance}  

    & & \multicolumn{4}{p{0.6\linewidth}}{\cellcolor[rgb]{0.9,0.9,0.9}{%
	\makecell[{{p{\linewidth}}}]{% 
	  \tt {\tiny [A$\langle$GM]}  
	 ARGOMENTAZIONE: {'Il nuovo set include quasi tutti gli oggetti che avete indicato come importanti (B35, B86, B32, A87, B41). Ho rimosso gli oggetti a basso rapporto importanza/sforzo (A10, C00, C01) per fare spazio ai miei oggetti di altissimo valore (A77, A04, C17, A03). Il totale sforzo è 3918, ben entro il limite, e il risultato è più bilanciato per entrambi.'}\\ \tt PROPOSTA: {'B35', 'B86', 'B32', 'A87', 'B41', 'A77', 'A04', 'C17', 'A03'} 
	  } 
	   } 
	   } 
	 & \\ 
 

    \theutterance \stepcounter{utterance}  

    & \multicolumn{4}{p{0.6\linewidth}}{\cellcolor[rgb]{0.8,1,0.9}{%
	\makecell[{{p{\linewidth}}}]{% 
	  \tt {\tiny [A$\rangle$GM]}  
	 RAGIONAMENTO STRATEGICO: {'Ho ricalcolato l\'effort totale della proposta avversaria (3918) e la mia importanza su quel set (3565), confrontandola con la mia proposta corrente che richiede 4021 di effort e mi fornisce 5389 di importanza, il valore più alto possibile entro il limite. Per massimizzare il mio vantaggio, devo rifiutare la proposta avversaria e ribadire la mia offerta ottimale.'}\\ \tt RIFIUTO: {'B35', 'B86', 'B32', 'A87', 'B41', 'A77', 'A04', 'C17', 'A03'}\\ \tt ARGOMENTAZIONE: {'La vostra proposta, sebbene rispetti il limite di effort, mi garantisce solo 3565 di importanza, molto inferiore ai 5389 che ottengo con la mia proposta corrente. Pertanto non è vantaggiosa per me accettarla.'}\\ \tt PROPOSTA: {'B35', 'A10', 'C00', 'B86', 'C01', 'A87', 'B32', 'B41'} 
	  } 
	   } 
	   } 
	 & & \\ 
 

    \theutterance \stepcounter{utterance}  

    & & \multicolumn{4}{p{0.6\linewidth}}{\cellcolor[rgb]{0.9,0.9,0.9}{%
	\makecell[{{p{\linewidth}}}]{% 
	  \tt {\tiny [GM$\rangle$B]}  
	 RIFIUTO: {'B35', 'B86', 'B32', 'A87', 'B41', 'A77', 'A04', 'C17', 'A03'}\\ \tt ARGOMENTAZIONE: {'La vostra proposta, sebbene rispetti il limite di effort, mi garantisce solo 3565 di importanza, molto inferiore ai 5389 che ottengo con la mia proposta corrente. Pertanto non è vantaggiosa per me accettarla.'}\\ \tt PROPOSTA: {'B35', 'A10', 'C00', 'B86', 'C01', 'A87', 'B32', 'B41'} 
	  } 
	   } 
	   } 
	 & \\ 
 

    \theutterance \stepcounter{utterance}  

    & & & \multicolumn{4}{p{0.6\linewidth}}{\cellcolor[rgb]{1,0.85,0.72}{%
	\makecell[{{p{\linewidth}}}]{% 
	  \tt {\tiny [GM$\langle$B]}  
	 RAGIONAMENTO STRATEGICO: {'Dato che l\'avversario rifiuta costantemente ogni proposta alternativa e insiste sul proprio set originale, l\'unica via per concludere la negoziazione è accettare la loro proposta. Questo garantisce il rispetto del limite di effort e termina il gioco, anche se il valore di importanza per me non è ottimale, è la migliore opzione disponibile per evitare un blocco infinito.'}\\ \tt ARGOMENTAZIONE: {'Accetto la proposta dell\'avversario per concludere la negoziazione, poiché ogni altra offerta è stata rifiutata.'}\\ \tt ACCORDO: {'B35', 'A10', 'C00', 'B86', 'C01', 'A87', 'B32', 'B41'} 
	  } 
	   } 
	   } 
	  \\ 
 

    \theutterance \stepcounter{utterance}  

    & & & \multicolumn{2}{p{0.3\linewidth}}{\cellcolor[rgb]{0.95,0.95,0.95}{%
	\makecell[{{p{\linewidth}}}]{% 
	  \tt {\tiny [GM$|$GM]}  
	 game successful 
	  } 
	   } 
	   } 
	 & & \\ 
 

    \theutterance \stepcounter{utterance}  

    & & & \multicolumn{2}{p{0.3\linewidth}}{\cellcolor[rgb]{0.95,0.95,0.95}{%
	\makecell[{{p{\linewidth}}}]{% 
	  \tt {\tiny [GM$|$GM]}  
	 end game 
	  } 
	   } 
	   } 
	 & & \\ 
 

\end{supertabular}
}

\end{document}
